\section{Discussion}
\label{sec:discussion}

\subsection{What a pass is evidence of}
\label{sec:discussion-evidence}

The four classes of reference support four different claims, and the results put a number on how much each supports. A printed answer shows that the engine implements the paper's formulas with the paper's conventions. The engine does so for \nPrintedAgree\ of \nPrintedN\ numbers at the strictest reading, \nAgreeOneUnit\ at one unit of the last digit, and it fails to agree in a way that the digits explain for \nMissTested\ of the \nMissTotal\ gaps and that recorded notes explain for \nMissRecorded\ more. This is what the reconciliation of a paper can show: \nMissNone\ gaps are unexplained, and they are bounded in size, the largest at \nMissNoneMax\ units of the last digit. It does not show that the paper is right, and the cases in which the engine and the paper differ by a few units may as well be the paper's rounding as the engine's choice of convention.

An independent route shows that the engine equals the mathematical definition to the tolerance, including the part of the definition that sits in a library. The measured gap to the definition is $\nGridJeffMax$ at the largest, more than two orders of magnitude below the tolerance, and the route flags every planted error from the tolerance upward. This is the strongest of the four kinds of evidence, and it is available only where a second computation of the definition exists, which for a tail probability it does and for a Monte Carlo estimator it does not. The reuse of the engine's own helper is the weakest reference and the most common one in a test suite; Proposition~\ref{prop:reuse} and the planted defects show that a test written that way proves the parameterisation and nothing else.

A closed form in a library shows that two code bases agree on one formula. The agreement here is to the last bit for the bond options and to $\nSwaptionPayer$ for the swaptions, and the grid shows where the relative gate loses its meaning: at a price of $\nClosedWorstShare\,\%$ of the notional. The gate is informative in proportion to the elasticity of the quantity to the parameter it checks, which for a deep in-the-money option is \nElasticityRatio\ times smaller than that of a swaption. A Monte Carlo reference shows that the difference of the estimates is small compared with their sampling error, which says little unless the rule that reads the difference is known to fail for a wrong engine. That the rule was never measured before is the main finding of this work: its size was above the nominal level for reasons that cancel, and its power was available only for defects of a few percent.

\subsection{Why the product's rule worked}
\label{sec:discussion-accident}

The product's rule accepted the single swap and the netting set and rejected the collateralised variants, which is the outcome that the exposure case deserves. The measurement of its size and power shows that it reached it by compensation. The standard error of each date's difference ignores the negative correlation of the antithetic pairs and is larger than the true error by a factor of \nOverstateLast\ at the last date, so the rule is less likely to reject at a date than its nominal bound says. The family of thirty-five comparisons at the bound of two would, with correct errors, reject two runs of the same model in \nSizePairNetting\,\% of the repetitions. The two effects meet at a size of \nSizeProdNetting\,\%, which is not nominal and is not alarming either. But the compensation depends on the case: a book whose exposures have weaker antithetic structure would have errors closer to correct and a rule that rejects much more often, and a profile with fewer dates would have a smaller family and a rule that is more conservative still. A reconciliation rule that happens to give the right outcome on one book has not shown that it would on another.

The repeated runs against ORE show the practical consequence: the rule passes the uncollateralised variants on most but not all seeds (\nReplNetPass\ of \nReplNetN\ and \nReplSinPass\ of \nReplSinN), so a single pair of runs that fails is not evidence of a defect, and one that passes is not evidence of equality. The refinement keeps the power of the product's rule against a small volatility defect with a size that is nominal, at the cost of two changes to the reconciliation code, one of them the estimator of the standard error. It is a measured improvement on this case and a proposal for others, not a result for books or engines that have not been run.

\subsection{The gaps and the negative results}
\label{sec:discussion-gaps}

Three of the six hypotheses are met in full (the independent routes, the Monte Carlo reference on the uncollateralised variants, and the listing of the collateralised gaps). Two are met with a qualification: one number of the wrong-way appendix is \nWWAtwoUnits\ units from its print, and one of the \nClosedCases\ closed-form cases is beyond the relative gate. One is not met as stated: the tally of printed numbers differs from the one the design document quotes, and the design document's counts are those of a tally that leaves out two tables. For the closed-form case and the tally the experiments identify the cause. For the printed constant they show that an independent bisection agrees with the engine, which places the difference in the printed value or its rounding without proving it. None of the three is a defect found in the engine.

The gaps against ORE are the one result that is a disagreement of two engines. The margining schedule explains part of the collateralised gaps and not all of them; the remainder is a property of the two engines on the seeds tried, since the 10-day variant fails the rule on \nReplCsaFail\ of \nReplCsaN\ new seeds. The initial-margin variant differs in the margin itself, ORE's at \nImRatio\ of the engine's, and the expected exposure after margin differs with it. The two regressions for the dynamic initial margin are different estimators, and which of them is nearer the margin that a clearing house would call is not something this comparison can say. Two negative results are worth stating. The CVA comparison of the product has no power against a small defect, because its error is the largest overstatement of the three; and the fixed relative tolerance, which looks like the simplest reconciliation rule, has either the size of a coin toss at one percent or no power below five.

\subsection{What a validation function should ask of a reconciliation}
\label{sec:discussion-ask}

The results suggest five questions. First, what do the reference and the engine share: a reference that reuses the engine's helper, or the library under it, cannot show a defect in them, and a report should say which function each reference calls. Second, what is the rule, fixed before the comparison, and what is its size and power, measured or derived; a rule with a random statistic and no stated error rates is a decoration. Third, what is the gap in a unit the reader can interpret, and what explains it, with the explanation tested where it can be. Fourth, what does the agreement support: an equivalence margin and a detectable gap say how much a pass is worth. Fifth, what is not claimed, with the reason, so that a pass on some units is not read as a pass on all. The last question is the one that the supervisory guidance on benchmarking already asks in other words \citep{fed2011sr117}, when it warns that a good match should be interpreted with caution; the protocol puts a number on the caution.
